\documentclass[10pt]{article}
\usepackage[T1]{fontenc}
\usepackage{lmodern}
\usepackage{amsmath,amssymb}
\usepackage[margin=0.72in]{geometry}
\usepackage[colorlinks=true,citecolor=blue,urlcolor=blue]{hyperref}

\newcommand{\C}{\mathbb C}
\newcommand{\A}{\mathcal A}

\begin{document}
\begin{center}
{\Large\bfseries Literature Status: The Factorial Group-Algebra
Completion\par}
\vspace{0.35em}
{Run \texttt{fa\_banach\_001}, agent \texttt{agent\_lane\_11}\par}
{\small 17 August 2026}
\end{center}
\vspace{0.4em}

\paragraph{Status.}
\textbf{literature\_already\_answered}.  A later primary paper by
Cahen--Gutt--Waldmann \cite{CGW} explicitly identifies and analyzes the
group-algebra completion left open by Beiser--Waldmann \cite{BW}.

\paragraph{Original questions.}
For a finitely generated infinite group $G$ with word length $L$, the source
rescales the group basis by
\[
 u_g=\frac{g}{(L(g)!)^\epsilon}\qquad(\epsilon>0)
\]
and defines a Fr\'echet topology by recursive seminorms.  Remark 3.14 asks
how the topology depends on the generators and on $\epsilon$, whether the
completion is exactly the formal series with subfactorial coefficient
growth, and for a useful tensor topology making the coproduct continuous.

\paragraph{Exact later identification.}
Section 5 of \cite{CGW} states that understanding this recursion was one
motivation for the later paper and reconstructs it with
$c(g)=(L(g)!)^\rho$.  Lemma 5.1 proves estimates in both directions, and
Proposition 5.2 proves equality of Fr\'echet algebras and of their locally
convex topologies:
\[
 \A_{L,\rho}(G)=\ell^1_{L,n!,\rho-}(G).
 \tag{1}\label{eq:id}
\]
The right side consists of canonical-basis series
$x=\sum_gx_g g$ such that
\[
 \sum_{g\in G}|x_g|(L(g)!)^R<\infty
 \qquad\text{for every }R<\rho.
 \tag{2}\label{eq:decay}
\]

This is precisely the requested subfactorial description in the source's
rescaled coordinates.  Indeed, if
$x=\sum_g b_gu_g$, then $x_g=b_g/(L(g)!)^\rho$, and
\eqref{eq:decay} becomes
\[
 \sum_{g\in G}\frac{|b_g|}{(L(g)!)^\delta}<\infty
 \qquad\text{for every }\delta>0,
 \tag{3}\label{eq:subfac}
\]
where it suffices to take $0<\delta<\rho$.  Thus the completion consists
exactly of subfactorially growing rescaled coefficients, with the
corresponding K\"othe topology.

\paragraph{Dependence on generators and parameter.}
Immediately after Proposition 5.2, the later paper states that
$\A_{L,\rho}(G)$ generally depends on the word length $L$; independence is
obtained only after taking the projective limit $\rho\to\infty$.  Formula
\eqref{eq:decay} also makes the parameter dependence explicit: increasing
$\rho$ requires faster factorial decay.  For an infinite group these
spaces are strictly nested, as seen on a geodesic ray $L(g_n)=n$ by taking
$x_{g_n}=(n!)^{-\alpha}$ and all other coefficients zero.

\newpage
\paragraph{Coproduct.}
For the weighted norm
\(
 \|x\|_{L,\sigma,R}=\sum_g|x_g|\sigma(L(g))^R,
\)
equation (4.7) of \cite{CGW} gives the exact relation
\[
 \|\Delta x\|_{R\widehat\otimes_\pi R}=\|x\|_{2R}.
 \tag{4}\label{eq:coproduct}
\]
This explains the fixed-window obstruction and the remedy: take the
projective limit over all $R\geq0$, so every target seminorm has the needed
$2R$ source seminorm.  Proposition 4.4 and Theorem 4.5 then give a
length-independent Fr\'echet--Hopf $^*$-algebra, functorial in finitely
generated groups.

\paragraph{Identification and scope.}
This is a direct later-paper answer, not a same-paper ask-and-answer passage.
The later paper cites \cite{BW}, reproduces its recursion, proves the exact
topological identification, and addresses the length and coproduct issues.
No new mathematical solution or priority claim is made here.

\begin{thebibliography}{9}
\small
\raggedright

\bibitem{BW}
S. Beiser and S. Waldmann,
\emph{Fr\'echet algebraic deformation quantization of the Poincar\'e disk},
J. Reine Angew. Math. \textbf{688} (2014), 147--207;
arXiv:1108.2004v3.

\bibitem{CGW}
M. Cahen, S. Gutt, and S. Waldmann,
\emph{Nuclear Group Algebras for Finitely Generated Groups},
Bull. Belg. Math. Soc. Simon Stevin \textbf{27} (2020), 567--594;
arXiv:1610.07746.

\end{thebibliography}

\end{document}
